\section{毕业设计题目}

\makeatletter
\@title % title macro of thesis, DO NOT MODIFY IT!
\makeatother

\section{主要任务与目标}

本次毕设课题主要任务为：基于Slidev框架设计并开发一款支持多人协作的智能幻灯片平台，整合AI辅助功能，实现云端资源管理、版本控制及Slidev格式深度兼容等核心功能。现将其应达到的目标与效果简要描述如下：
\begin{enumerate}
    \item 支持多人实时协同编辑、角色权限分配及版本控制与回溯功能；
    \item 实现云端统一存储图片、模板等资源，支持团队共享与快速调用；
    \item 提供AI辅助内容生成、格式优化等智能制作功能；
    \item 深度兼容Slidev原生功能，包括Latex公式、代码高亮、Mermaid图表等；
\end{enumerate}

\section{主要内容与基本要求}

本平台开发需遵循软件工程规范，依次完成需求分析、架构设计、数据库设计及界面布局规划，具体内容与要求如下：
\begin{enumerate}
    \item 主要内容：
    \begin{enumerate}
        \item 支持多人实时协同编辑、角色权限分配及版本控制与回溯功能；
        \item 实现云端统一存储图片、模板等资源，支持团队共享与快速调用；
        \item 提供AI辅助内容生成、格式优化等智能制作功能；
        \item 深度兼容Slidev原生功能，包括Latex公式、代码高亮、Mermaid图表等；
    \end{enumerate}
    \item 基本要求：
    \begin{enumerate}
        \item 多人协作管理模块：实现多人实时协同编辑功能，采用基于角色的权限分配系统\cite{zhang2013rbac,wu2002rbacweb,kuhlmann2013uml}，参考UML权限设计规范\cite{kuhlmann2013uml}及软件重构可行性经验\cite{li2011reengineering}，支持基础角色权限设置；
        \item 云端资源管理模块：搭建云端资源存储与共享功能，参考文档共享系统方案\cite{li2015filenet}及软件项目文档管理系统的资源组织形式\cite{zhang2024javaweb}，实现资源的基础检索与调用；
        \item AI 辅助功能模块：集成基础的 AI 内容生成、格式优化工具，参考插件化设计思路\cite{cai2022microfrontend}保障功能扩展性，提升幻灯片制作便捷性；
        \item 按照基础软件工程流程完成编码实现、简单模块测试与系统优化，参考云资源管理中的应用实践\cite{wei2025rbaccloud}及WEB文档协作系统设计思路\cite{tan2012web,li2020cscw}，确保平台稳定运行，满足团队协作制作幻灯片的核心需求。
    \end{enumerate}
\end{enumerate}

\section{毕业设计进度安排}
\begin{table}[!htbp]
	\centering
	\vspace{-2ex}
	\small% fontsize
	\begin{tabular}{cp{25em}p{10.5em}}
		\toprule
		\textbf{序号} & \multicolumn{1}{c}{\textbf{进度}} & \multicolumn{1}{c}{\textbf{起止时间}} \\
		\midrule
		1           & 任务下达                          & 2025.11.27 $\sim$ 2025.12.03      \\
		2           & 开题报告、文献综述等文档撰写      & 2025.12.06 $\sim$ 2025.12.27      \\
		3           & 完成开题答辩                      & 2025.12.27 $\sim$ 2026.01.02      \\
		4           & 整体设计与初步开发                & 2026.01.02 $\sim$ 2026.04.01      \\
		5           & 完成中期检查                      & 2026.04.01 $\sim$ 2026.04.15      \\
		6           & 项目开发收尾与毕业论文撰写        & 2026.04.15 $\sim$ 2026.05.01      \\
		7           & 答辩准备与成果整理                & 2026.05.01 $\sim$ 2026.05.17      \\
		8           & 完成答辩与材料归档                & 2026.05.17 $\sim$ 2026.05.20      \\
		\bottomrule
	\end{tabular}
	\vspace{-2ex}
\end{table}

\section{主要参考文献}